\clearpage\section*{EV4-A -- Préparation EAM}\label{EV4-A}
\addcontentsline{toc}{section}{EV4-A -- Préparation EAM}
\textbf{Sujet réservé au professeur}. Ensemble du programme.\par Durée : 2 h. Sans calculatrice. Note sur 20. Entraînements de mai et juin ; EAM de zone à confirmer.
\subsection*{Partie 1 -- Automatismes (6 points)}
Douze questions, 0,5 point chacune. Une seule réponse correcte ; aucune justification demandée. Une réponse incorrecte ou absente rapporte 0, sans retrait de points.
\begin{enumerate}
\item Calculer $2+2\times3$.\par A. $11$\quad B. $7$\quad C. $8$\quad D. $12$
\item Calculer $\frac12+\frac13$.\par A. $\frac{5}{3}$\quad B. $\frac{1}{6}$\quad C. $\frac{5}{6}$\quad D. $\frac{2}{5}$
\item Calculer $\dfrac{2^2}{2^{0}}$.\par A. $2$\quad B. $1$\quad C. $7$\quad D. $4$
\item Une valeur de 50 augmente de 20 \%. Valeur finale ?\par A. $70$\quad B. $40$\quad C. $60$\quad D. $10$
\item Une valeur de 80 est obtenue après une baisse de 20 \%. Valeur initiale ?\par A. $100$\quad B. $60$\quad C. $96$\quad D. $64$
\item Hausse de 25 \% puis baisse de 20 \% : coefficient global ?\par A. $1$\quad B. $\frac{4}{5}$\quad C. $\frac{21}{20}$\quad D. $\frac{9}{10}$
\item La solution de $2x+2=0$ vaut\par A. $1$\quad B. $4$\quad C. $-1$\quad D. $-4$
\item La pente entre $(0;2)$ et $(2;8)$ vaut\par A. $-3$\quad B. $3$\quad C. $6$\quad D. $2$
\item La moyenne de $2, 4, 6$ vaut\par A. $4$\quad B. $5$\quad C. $2$\quad D. $12$
\item Un secteur de $90^\circ$ représente quelle fraction du total ?\par A. $\frac{1}{4}$\quad B. $\frac{1}{3}$\quad C. $\frac{3}{4}$\quad D. $\frac{1}{2}$
\item Une probabilité vaut $\dfrac{2}{5}$. Son complément vaut\par A. $\frac{5}{2}$\quad B. $\frac{3}{5}$\quad C. $\frac{2}{5}$\quad D. $\frac{1}{5}$
\item Pour $f(x)=x^2+2$, $f(-2)$ vaut\par A. $6$\quad B. $-2$\quad C. $4$\quad D. $8$
\end{enumerate}
\clearpage\subsection*{Partie 2 -- Exercices indépendants (14 points)}
Toute démarche doit être justifiée. Les résultats exacts sont attendus sauf indication contraire.
\Needspace{12\baselineskip}
\subsection*{Exercice 1 -- Deux variables liées (5 points)}
La variable aléatoire $X$ prend les valeurs 1, 2 et 3, chacune avec probabilité $\dfrac13$. On définit sur la même expérience la variable $Y=4-X$. On souhaite vérifier si une règle valable pour l’espérance reste vraie pour la variance.
\begin{enumerate}
\item Calculer $E(X)$ et $E(Y)$. \hfill(1 pt)
\item Calculer $V(X)$, $V(Y)$ et $V(X+Y)$. \hfill(2 pts)
\item Utiliser ces résultats pour réfuter l’affirmation : « La variance d’une somme est toujours égale à la somme des variances. » \hfill(2 pts)
\end{enumerate}
\Needspace{12\baselineskip}
\subsection*{Exercice 2 -- Comparer courbe et droite (5 points)}
Les fonctions $f(x)=e^{-x}$ et $g(x)=1-x$ sont définies sur $\R$. On pose $h=f-g$. La comparaison de leurs courbes permettra ensuite de contrôler le sens d’une approximation.
\begin{enumerate}
\item Dériver leur différence. \hfill(1 pt)
\item Prouver $f(x)\geq g(x)$ pour tout réel x. \hfill(2 pts)
\item Écrire l’approximation affine de $e^{-0{,}02}$ et préciser si elle est par défaut ou par excès. \hfill(2 pts)
\end{enumerate}
\Needspace{12\baselineskip}
\subsection*{Exercice 3 -- Vecteur normal et rayon (4 points)}
Dans un repère orthonormé, on considère le point $I(1;1)$ et la droite $d$ d’équation $3x+4y-17=0$. On note $H$ le projeté orthogonal de $I$ sur $d$. On cherche le cercle de centre $I$ tangent à $d$.
\begin{enumerate}
\item Calculer la distance de I à d. \hfill(1 pt)
\item Calculer le projeté H. \hfill(1 pt)
\item Donner l’équation du cercle de centre I tangent à d et vérifier que H appartient au cercle. \hfill(2 pts)
\end{enumerate}
